\section{Related Work}

\noindent{\bf Denoising Diffusion Models for Human Animation.}
Denoising diffusion models~\citep{sohl2015deep,song2020denoising,ho2020denoising} synthesize data by gradually denoising samples drawn from a noise distribution, effectively reversing a stochastic diffusion process. Recently, these models have achieved notable success in human motion generation, producing realistic and diverse animations~\citep{barquero2022belfusion,tevet2023human,zhang2022motiondiffuse,raab2023single,zhang2023t2m,yonatan2023,zhang2023tedi}. A prominent example, MDM~\citep{tevet2023human}, leverages transformer architectures to predict clean human motion trajectories from noisy inputs during the denoising stages. To generate conditional motion sequences, such diffusion methods typically incorporate external conditions into the diffusion process, including textual descriptions~\citep{petrovich2023tmr,guo2022tm2t,petrovich22temos,zhang2022motiondiffuse,zhang2023motiongpt,tevet2022motionclip,barquero2024seamless}. Further advancements have extended diffusion-based approaches to more conditioning scenarios, including guiding animations along specific trajectories~\citep{karunratanakul2023gmd,rempe2023trace,xie2023omnicontrol}.

\noindent \textbf{Denoising Diffusion Models for Interaction Animation.}
Recent advancements in interaction animation have shown significant progress in diverse scenarios, including human-human interactions~\citep{liang2023intergen,CMU-Mocap,singleshotmultiperson2018,xu2023inter, xu2023stochastic} and human interactions within static scenes~\citep{hassan_samp_2021,cao2020long,hassan2019resolving}. Early human-object interaction (HOI) animation methods typically leveraged kinematic models combined with conventional generative frameworks. These earlier approaches primarily addressed simplified scenarios, either involving static or small objects~\citep{xie2022chore,zhang2020perceiving,wang2022reconstructing,petrov2023object,taheri2022goal,wu2022saga,kulkarni2023nifty,zhang2022couch}, or focusing exclusively on hand-object interactions~\citep{li2023task,ye2023affordance,zheng2023cams,zhou2022toch,zhang2024manidext,zhang2023artigrasp}. Consequently, they struggled to effectively capture complex, dynamic, whole-body interactions. Recently, denoising diffusion models have emerged as a powerful paradigm capable of modeling sophisticated interactions involving dynamically moving objects~\citep{peng2023hoi,diller2023cg,li2023controllable,wu2024thor,wu2024human,song2024hoianimator,xu2024interdreamer,zhang2024hoi}. 
Concurrently, new datasets have expanded the range of possible interactions -- from low-dynamic manipulations~\citep{taheri2020grab} to highly dynamic scenarios engaging multiple body parts simultaneously~\citep{bhatnagar22behave,jiang2022chairs,huang2022intercap,zhang2023neuraldome,fan2023arctic,li2023object,zhao2023im,kim2024parahome,jiang2024scaling,yang2024f}. We use the most comprehensive InterAct dataset~\citep{xu2025interact} for evaluation.



\noindent \textbf{Guidance with Denoising Diffusion Models.}
Diffusion models can be steered by additional guidance signals to better satisfy conditioning constraints. In image and text generation, for instance, classifier guidance~\citep{dhariwal2021diffusion} uses an external classifier’s gradients to direct the denoising process toward desired outcomes, whereas classifier-free guidance~\citep{ho2022classifier} foregoes a separate classifier by leveraging the model’s own conditional and unconditional predictions for guidance. Inspired by such strategies, recent works in 3D human-object interaction (HOI) generation have explored explicit contact-aware or auxiliary-guided approaches to ensure realism. Without explicit constraints, diffusion-based HOI models often produce unrealistic artifacts, e.g., floating objects or missing contact points. To mitigate this, InterDiff~\citep{xu2023interdiff}, for example, employs a kinematics-informed predictor that iteratively refines the diffusion output with corrections, improving human-object prediction accuracy. Similarly, HOI-Diff~\citep{peng2023hoi} integrates an auxiliary affordance module to steer the model toward consistent contact and affordance cues throughout generation. These approaches demonstrably boost realism but introduce additional complexity by relying on extra networks or hand-crafted constraints, which can hinder generalization. CG-HOI~\citep{diller2023cg} cast interaction synthesis as a multi-task learning problem: an auxiliary contact prediction task is learned jointly to guide the motion generation, thereby avoiding training separate guidance models. However, it still imposes predefined relationships between predicted contact and object movements, using weighted combinations that embed human-designed assumptions. In contrast to these externally guided methods, we propose a guidance strategy without any assumptions on human-object contact, where the proposed approach gradually applies guidance according to a temporal schedule for different components during denoising, promoting a new principled manner and and potentially inspiring future applications in other domains.
